% Options for packages loaded elsewhere
\PassOptionsToPackage{unicode}{hyperref}
\PassOptionsToPackage{hyphens}{url}
\PassOptionsToPackage{dvipsnames,svgnames,x11names}{xcolor}
%
\documentclass[
  11pt,
]{article}
\usepackage{amsmath,amssymb}
\usepackage{iftex}
\ifPDFTeX
  \usepackage[T1]{fontenc}
  \usepackage[utf8]{inputenc}
  \usepackage{textcomp} % provide euro and other symbols
\else % if luatex or xetex
  \usepackage{unicode-math} % this also loads fontspec
  \defaultfontfeatures{Scale=MatchLowercase}
  \defaultfontfeatures[\rmfamily]{Ligatures=TeX,Scale=1}
\fi
\usepackage{lmodern}
\ifPDFTeX\else
  % xetex/luatex font selection
\fi
% Use upquote if available, for straight quotes in verbatim environments
\IfFileExists{upquote.sty}{\usepackage{upquote}}{}
\IfFileExists{microtype.sty}{% use microtype if available
  \usepackage[]{microtype}
  \UseMicrotypeSet[protrusion]{basicmath} % disable protrusion for tt fonts
}{}
\makeatletter
\@ifundefined{KOMAClassName}{% if non-KOMA class
  \IfFileExists{parskip.sty}{%
    \usepackage{parskip}
  }{% else
    \setlength{\parindent}{0pt}
    \setlength{\parskip}{6pt plus 2pt minus 1pt}}
}{% if KOMA class
  \KOMAoptions{parskip=half}}
\makeatother
\usepackage{xcolor}
\usepackage[margin=1in]{geometry}
\usepackage{longtable,booktabs,array}
\usepackage{calc} % for calculating minipage widths
% Correct order of tables after \paragraph or \subparagraph
\usepackage{etoolbox}
\makeatletter
\patchcmd\longtable{\par}{\if@noskipsec\mbox{}\fi\par}{}{}
\makeatother
% Allow footnotes in longtable head/foot
\IfFileExists{footnotehyper.sty}{\usepackage{footnotehyper}}{\usepackage{footnote}}
\makesavenoteenv{longtable}
\setlength{\emergencystretch}{3em} % prevent overfull lines
\providecommand{\tightlist}{%
  \setlength{\itemsep}{0pt}\setlength{\parskip}{0pt}}
\setcounter{secnumdepth}{-\maxdimen} % remove section numbering
\DeclareUnicodeCharacter{220E}{\ensuremath{\square}}
\ifLuaTeX
  \usepackage{selnolig}  % disable illegal ligatures
\fi
\IfFileExists{bookmark.sty}{\usepackage{bookmark}}{\usepackage{hyperref}}
\IfFileExists{xurl.sty}{\usepackage{xurl}}{} % add URL line breaks if available
\urlstyle{same}
\hypersetup{
  pdftitle={Conditional Davenport Corridors and Saturated Obstructions in C5 Cubed},
  pdfauthor={Sze Chun Yiu},
  pdfsubject={Conditional Davenport corridors and saturated obstructions in C5 cubed},
  pdfkeywords={generalized Davenport constants, zero-sum sequences, elementary abelian groups},
  colorlinks=true,
  linkcolor={Maroon},
  filecolor={Maroon},
  citecolor={Blue},
  urlcolor={Blue},
  pdfcreator={LaTeX via pandoc}}

\title{Conditional Davenport Corridors and Saturated Obstructions in \(C_5^3\)}
\author{Sze Chun Yiu\\\texttt{sze-chun.yiu@fysik.su.se}}
\date{25 August 2026}

\begin{document}
\maketitle

\begin{abstract}
Let \(D_k(G)\) be the least length that forces \(k\) pairwise disjoint
nonempty zero-sum subsequences. Freeze and Schmid proved the recurrence
and the lower line \(D_k(C_5^3)\ge5k+10\); Fan et al.~record
\(s_{\le5}(C_5^3)=33\) and a support-eight inverse property attributed
to Gao et al.~We derive \(D_2(C_5^3)=20\) from Freeze--Schmid's lower
bound and Zhao's short-subsequence lemma. Assuming only
\(D_3(C_5^3)=25\) and \(s_{\le6}(C_5^3)=24\), we obtain
\(5k+10\le D_k(C_5^3)\le5k+11\) for every \(k\ge4\).

If \(D_4(C_5^3)=30\), then the lower line is exact for every \(k\ge2\).
To analyze the remaining alternative, we study a hypothetical total-zero
sequence of length \(31\) with no nonempty zero-sum subsequence of
length at most five. For saturated short-free sequences in odd-prime
elementary abelian groups, we prove a defect certificate excluding
multiplicity \(p-2\). For \(p=5\), a saturated obstruction has only
multiplicities \(1,2,4\). If \(s\) is support size and \(c_i\) counts
multiplicity \(i\), then \(c_2=31-s-3c_4\) and \(c_1=2s-31+2c_4\).

The subsequence of points of multiplicity at least two has length
\(62-2s-2c_4\). Using the short-zero-sum threshold \(\eta(C_5^2)=13\),
we prove that this stratum spans rank three whenever \(s+c_4\le24\). We
also prove an exact identity coupling support deficit, repetitions
inside zero-sum atoms, and overlaps between atom supports. Two exact
implementations found no obstruction through support ten within the
supplied computational model. Because this computation has not been
independently replicated, we report it only as bounded computational
evidence.

Under H1 and H2, these results leave \(D_4(C_5^3)\in\{30,31\}\).

\end{abstract}

\noindent\textbf{Keywords:} generalized Davenport constants; zero-sum sequences;
elementary abelian groups; short zero sums; inverse zero-sum problems

\hypertarget{introduction}{%
\section{1. Introduction}\label{introduction}}

Generalized Davenport constants measure the sequence length required to
force several pairwise disjoint zero sums. Rank-two groups exhibit
strong eventual structure, whereas rank-three groups can depart from the
expected arithmetic progression.

Under the declared inputs below, the values \(D_2=20\) and \(D_3=25\),
together with short-zero-sum thresholds, reduce the next value to

\[
D_4(C_5^3)\in\{30,31\}.
\]

This one-bit ambiguity controls the tail: the lower candidate \(D_4=30\)
forces \(D_k=5k+10\) for every \(k\ge2\). The upper candidate would
yield a rigid length-31 total-zero sequence with no zero sum of length
at most five.

We combine the recurrence with symbolic analysis of that hypothetical
obstruction. The purpose is twofold: to obtain a conditional width-one
theorem for the full tail and to identify exactly where existing rank
reductions cease to be automatic. Computation is kept in a separate
evidence class.

Our contributions are:

\begin{enumerate}
\def\labelenumi{\arabic{enumi}.}
\tightlist
\item
  the conditional width-one corridor for every \(k\ge4\);
\item
  the conditional exact tail if \(D_4=30\);
\item
  an odd-prime saturation-defect lemma excluding multiplicity \(p-2\);
\item
  bounded computational evidence through support ten;
\item
  an exact multiplicity parameterization for every remaining support
  stratum;
\item
  the rank-forcing phase \(s+c_4\le24\) for the high-multiplicity
  stratum; and
\item
  an identity coupling support deficit to internal atom repetition and
  cross-atom overlap.
\end{enumerate}

\hypertarget{definitions-and-declared-inputs}{%
\section{2. Definitions and declared
inputs}\label{definitions-and-declared-inputs}}

For a finite abelian group \(G\), \(D_k(G)\) is the least integer \(n\)
such that every sequence of length at least \(n\) over \(G\) contains
\(k\) pairwise disjoint nonempty zero-sum subsequences. Let
\(s_{\le\ell}(G)\) denote the least length forcing a nonempty zero-sum
subsequence of length at most \(\ell\). For an odd prime \(p\), a
sequence is \emph{\(p\)-short-free} if it has no nonempty zero-sum
subsequence of length at most \(p\). It is \emph{saturated} relative to
this property if appending any group element destroys it. Thus
\emph{5-short-free} and \emph{saturated} below are the specializations
at \(p=5\).

We use the generalized-Davenport recurrence of Freeze and Schmid {[}1,
Proposition 3.1{]}

\[
D_{k+1}(G)
\le\max\{D_k(G)+\ell, \ s_{\le\ell}(G)-1\},
\]

and their lower bound \(D_k(C_5^3)\ge5k+10\) {[}1, Theorem 4.1{]}. Fan
et al. record \(s_{\le5}(C_5^3)=\eta(C_5^3)=33\) and Property C for
\(C_5^3\) {[}2, Lemma 4.2(5),(8){]}, attributing the original results to
Gao et al.~{[}6{]}. By the definition of Property C, every length-32
5-short-free sequence has exactly eight distinct support points. We also
use the rank-two identity \(\eta(C_m^2)=3m-2\), hence \(\eta(C_5^2)=13\)
{[}2, p.~1{]}.

For clarity, the two assumptions not established in the cited literature
are

\[
\text{(H1)}\quad D_3(C_5^3)=25,
\qquad
\text{(H2)}\quad s_{\le6}(C_5^3)=24.
\]

Every corridor or obstruction statement below that uses H1 or H2 says so
explicitly. Neither assumption is presented as a theorem of this paper.

\textbf{Proposition 1.} \(D_2(C_5^3)=20\).

\textbf{Proof.} Freeze and Schmid's lower bound gives
\(D_2(C_5^3)\ge20\). For the reverse inequality, Olson's formula for
finite abelian \(p\)-groups gives \(D(C_5^3)=1+3(5-1)=13\) {[}7{]}. Let
\(S\) be an arbitrary sequence of length 20 and append \(g=-\sigma(S)\),
producing the zero-sum sequence \(T=Sg\) of length 21. Apply Zhao's
Lemma 4.4 {[}3{]} with \(p=5\), \(D(C_5^3)=13\), \(|T|=21\), \(k=8\),
and \(i=2\). Its coefficient is
\(\binom{13}{6}+\binom{14}{7}=5148\equiv3\pmod5\), so \(T\) has a
nonempty zero-sum subsequence of length at most seven. Its zero-sum
complement has length at least 14 and therefore contains a nonempty
proper zero-sum subsequence because \(D(C_5^3)=13\); the remaining
complement is zero-sum as well. Thus \(T\) contains three pairwise
disjoint nonempty zero-sum subsequences. At most one contains the
appended term \(g\), so the other two are disjoint zero-sum subsequences
of the arbitrary sequence \(S\). Therefore \(D_2(C_5^3)\le20\). ∎

\hypertarget{a-width-one-tail}{%
\section{3. A width-one tail}\label{a-width-one-tail}}

Taking \(\ell=6\) gives

\[
D_4\le\max\{25+6,24-1\}=31.
\]

Together with the lower bound, this proves \(D_4\in\{30,31\}\).

\textbf{Theorem 2 (conditional corridor).} Under H1 and H2, for every
\(k\ge4\),

\[
5k+10\le D_k(C_5^3)\le5k+11.
\]

\textbf{Proof.} The base case is \(D_4\le31\). Apply the recurrence with
\(\ell=5\) and \(s_{\le5}-1=32\). If \(D_k\le5k+11\), then both terms in
the maximum are at most \(5(k+1)+11\). Freeze and Schmid's lower bound
supplies the lower line. ∎

\textbf{Theorem 3 (conditional tail).} Under H1 and H2, if
\(D_4(C_5^3)=30\), then

\[
D_k(C_5^3)=5k+10
\qquad(k\ge2).
\]

\textbf{Proof.} The statement holds for \(k=2,3,4\). Applying the same
recurrence with \(\ell=5\) propagates the upper bound \(5(k+1)+10\),
which matches the published lower bound. ∎

There is no converse here: the assumption \(D_4=31\) does not by itself
prove that later terms lie on the upper line.

\hypertarget{the-hypothetical-length-31-obstruction}{%
\section{4. The hypothetical length-31
obstruction}\label{the-hypothetical-length-31-obstruction}}

Assume H1 and H2, and suppose \(D_4(C_5^3)=31\). The extremal
characterization of Freeze and Schmid provides a total-zero sequence
\(S\) of length 31 whose maximum factorization length is four. If \(M\)
is the minimum length of an atom dividing \(S\), their recurrence gives
\(31=D_4\le D_3+M=25+M\), so \(M\ge6\). On the other hand,
\(s_{\le6}(C_5^3)=24\) forces every 31-term sequence to contain a
zero-sum subsequence of length at most six; an inclusion-minimal such
subsequence is an atom. Hence \(M=6\), and \(S\) contains no zero-sum
subsequence of length at most five. Thus the upper candidate yields a
sequence satisfying

\[
|S|=31, \qquad \sigma(S)=0,
\]

with no nonempty zero-sum subsequence of length at most five. Every
element has multiplicity at most four. If its support exceeds eight, it
is saturated: otherwise a short-free one-term extension would have
length 32, whereas Property C {[}2{]} forces every length-32
5-short-free sequence to have exactly eight support points.

\textbf{Theorem 4 (saturation defect).} Let \(p\) be odd and let \(S\)
be a saturated \(p\)-short-free sequence over an elementary abelian
exponent-\(p\) group. If a nonzero element \(x\) has multiplicity
\(m<p\), then there is a subsequence \(R\) such that

\[
|R|\le p-1-m,
\qquad x\notin\operatorname{supp}(R),
\qquad \sigma(R)=-(m+1)x.
\]

\textbf{Proof.} Append one further copy of \(x\). Saturation supplies a
zero-sum subsequence of length at most \(p\) that uses the appended
occurrence. It must also use every original copy of \(x\); otherwise
replacing the appended copy by an unused original occurrence would give
a short zero sum already contained in \(S\). Removing the \(m+1\) copies
of \(x\) leaves the required \(R\). ∎

\textbf{Corollary 5.} Multiplicity \(p-2\) is impossible.

Indeed, the defect subsequence would have length at most one and sum
\(x\), forcing a forbidden additional copy of \(x\). For \(p=5\),
multiplicity three is absent. Singleton and double points also acquire
defect certificates of length at most three and two, respectively.

\hypertarget{exact-multiplicity-grammar}{%
\section{5. Exact multiplicity
grammar}\label{exact-multiplicity-grammar}}

For the saturated case \(s\ge9\), Corollary 5 excludes multiplicity
three. Let \(s=|\operatorname{supp}(S)|\), and let \(c_1, c_2, c_4\)
count points of multiplicity one, two, and four. Then

\[
c_1+c_2+c_4=s,
\qquad
c_1+2c_2+4c_4=31.
\]

Writing \(c_4=j\) gives

\[
c_2=31-s-3j,
\qquad
c_1=2s-31+2j.
\]

Thus \(c_1\) is odd, and nonnegativity determines the complete
admissible range for every support size. For example, the support-23
patterns are

\[
1^{15}2^8, \qquad 1^{17}2^5 4, \qquad 1^{19}2^2 4^2,
\]

and the support-24 patterns are

\[
1^{17}2^7, \qquad 1^{19}2^4 4, \qquad 1^{21}2 4^2.
\]

\hypertarget{the-rank-forcing-phase}{%
\section{6. The rank-forcing phase}\label{the-rank-forcing-phase}}

Let \(H\) be the subsequence consisting of all points of multiplicity
two or four. Its length is

\[
|H|=2c_2+4c_4=62-2s-2c_4.
\]

If \(H\) were contained in a subgroup of rank at most two, it would be a
5-short-free sequence over \(C_5^2\). The value \(\eta(C_5^2)=13\) {[}2,
p.~1{]} forces a short zero sum once \(|H|\ge13\). Since \(|H|\) is
even, the contradiction begins at \(|H|\ge14\).

\textbf{Theorem 6 (rank-forcing phase).} For a saturated obstruction
with \(s\ge9\), the high-multiplicity stratum spans rank three whenever

\[
s+c_4\le24.
\]

\textbf{Proof.} The condition \(|H|\ge14\) is equivalent to
\(62-2s-2c_4\ge14\), hence to \(s+c_4\le24\). ∎

The theorem gives the following exact boundary:

\begin{longtable}[]{@{}rrl@{}}
\toprule\noalign{}
Support & \(c_4\) & Consequence from Theorem 6 \\
\midrule\noalign{}
\endhead
\bottomrule\noalign{}
\endlastfoot
23 & 0 or 1 & rank three forced \\
23 & 2 & threshold is silent \\
24 & 0 & rank three forced \\
24 & 1 or 2 & threshold is silent \\
at least 25 & any admissible value & threshold is silent \\
\end{longtable}

Silence is not a converse. Outside the phase, rank three may follow from
a different argument.

\hypertarget{bounded-low-support-exclusion}{%
\section{7. Bounded low-support
exclusion}\label{bounded-low-support-exclusion}}

Symbolic reductions combined with two exact state representations
exclude every support size through ten.

\textbf{Computational result.} In the supplied exact search model, no
length-31 total-zero 5-short-free sequence over \(C_5^3\) exists with
support at most ten.

\textbf{Exact search description.} Multiplicity at most four first
excludes supports at most seven. At support eight, the only length
pattern is \(4^7 3\). Adding the missing fourth copy of the triple point
preserves short-freeness; total sum then forces that point to equal the
negative support sum. After normalizing a basis by \(\mathrm{GL}(3,5)\),
two exact subset-sum representations enumerate all 564 normalized
supports and find none satisfying this condition.

For supports nine and ten, saturation and Corollary 5 exclude
multiplicity three. The support-nine equations leave only \(4^7 2 1\); a
canonical search with the last point forced by total sum visits
6,537,270 states in each of two distinct exact-weight subset-sum engines
and finds no solution. At support ten, the only patterns are \(1^3 4^7\)
and \(1,2^3 4^6\). Four multiplicity-four points must span rank three
because 16 terms in a rank-two subgroup would contradict
\(\eta(C_5^2)=13\) {[}2, p.~1{]}. Normalizing an independent triple to
the standard basis makes the remaining enumeration complete. The two
engines agree exactly: they visit respectively 210,700 states with 3,558
leaves and 272,119 states with no leaves, finding zero solutions in both
patterns.

Both engines prune a partial candidate exactly when adding a term
creates a zero-sum subsequence of one of the lengths one through five;
one stores explicit byte reachability by weight and group sum, while the
other stores translation masks in a 128-bit representation. Their
source, build instructions, expected rows, and machine-readable results
are included as ancillary files. Thus every support stratum through ten
is exhausted within the declared model. Because a clean-room external
replay has not been completed, this result is not used in any analytic
proof.

\hypertarget{atom-repetition-and-overlap}{%
\section{8. Atom repetition and
overlap}\label{atom-repetition-and-overlap}}

Let \(S=U_1\cdots U_r\) be an occurrence-disjoint atom factorization.
For each group element \(g\), let \(r_g\) be the number of atom supports
containing \(g\), and define the internal deficit

\[
\delta_i=|U_i|-|\operatorname{supp}(U_i)|.
\]

\textbf{Theorem 7 (repetition-overlap identity).} Every
occurrence-disjoint factorization satisfies

\[
|S|-|\operatorname{supp}(S)|
=\sum_i\delta_i
+\sum_{g\in\operatorname{supp}(S)}(r_g-1).
\]

\textbf{Proof.} Summing the atom support sizes counts every global
support point \(r_g\) times. Hence

\[
\sum_i\delta_i=|S|-\sum_g r_g,
\]

while cross-atom overlap equals

\[
\sum_g r_g-|\operatorname{supp}(S)|.
\]

Adding the two identities proves the claim. ∎

Thus any independently established lower support bound \(s_0\) gives a
total internal-repetition plus cross-overlap budget of at most
\(31-s_0\). The two packaged support-ten implementations suggest
\(s_0=11\) and hence budget 20, but that numerical consequence remains
bounded computational evidence and is not used in any proof.

\hypertarget{the-residual-obstruction-phase}{%
\section{9. The residual obstruction
phase}\label{the-residual-obstruction-phase}}

A decisive next analysis should impose simultaneously:

\begin{enumerate}
\def\labelenumi{\arabic{enumi}.}
\tightlist
\item
  the multiplicity pattern from Section 5;
\item
  the rank status from Theorem 6;
\item
  the saturation-defect certificate at every singleton and double point;
\item
  the repetition-overlap identity; and
\item
  total sum zero with exclusion of zero sums of lengths one through
  five.
\end{enumerate}

This intersection is substantially smaller than an unstructured support
search. A surviving sequence would require independent factorization
verification before establishing \(D_4=31\). A complete, independently
replayed infeasibility proof for the length-31 target would establish
the lower candidate \(D_4=30\).

\hypertarget{relation-to-prior-work}{%
\section{10. Relation to prior work}\label{relation-to-prior-work}}

The generalized-Davenport recurrence, lower-bound framework, extremal
factorization language, and eventual-linearity context are established
results {[}1{]}. The value \(D_2(C_5^3)=20\) is a direct corollary of
the published lower bound and Zhao's short-subsequence lemma {[}3{]},
not a novelty claim here. The \(C_0\) framework and short-zero-sum
localization are also prior work {[}2{]}, as are the inverse zero-sum
program {[}4{]} and rank-two inverse results {[}5{]}.

The residual contribution is the width-one synthesis for \(C_5^3\), the
saturation-defect specialization, the exact multiplicity grammar, the
rank-forcing phase, and the coupling of support deficit with atom
repetition and overlap. These are structural compression results; they
do not decide the final extremal bit.

\hypertarget{reproducibility-and-limitations}{%
\section{11. Reproducibility and
limitations}\label{reproducibility-and-limitations}}

The conditional corridor, defect lemma, and support-ten exclusion are
separately bound to exact records. A separate implementation enumerates
all multiplicity patterns for supports 8-31, checks the rank-phase
equivalence, reproduces the support-23 and support-24 tables, and
validates the atom identity on independently generated factorizations.
Only the displayed analytic proofs carry all-parameter authority.

The exact value of \(D_4(C_5^3)\) remains open, as does the associated
length-31 extremal-spectrum statement. The rank phase is one-way, and
the atom identity is a compression principle rather than a
contradiction. No additional implementation can substitute for an
independently auditable final obstruction proof.

\hypertarget{conclusion}{%
\section{12. Conclusion}\label{conclusion}}

Conditional on H1 and H2, the generalized Davenport constants of
\(C_5^3\) remain within one unit of the line \(5k+10\) for every
\(k\ge4\), and the lower choice at \(k=4\) would fix the entire tail.
The hypothetical upper-line obstruction has substantially more structure
than this corridor alone reveals: its multiplicities have an exact
grammar, its high-multiplicity stratum enters a sharp rank-forcing
phase, and its atom factorization obeys a global repetition-overlap
budget.

These results explain where the residual difficulty begins and define a
smaller intersection problem for the unresolved extremal bit. They do
not claim to have resolved it.

\hypertarget{tool-use-disclosure}{%
\section{Tool-use disclosure}\label{tool-use-disclosure}}

A generative language model assisted manuscript organization, language
revision, and submission-package preparation. The listed author remains
responsible for the mathematical statements, proofs, references,
executable claims, and final text.

\hypertarget{data-and-code-availability}{%
\section{Data and code availability}\label{data-and-code-availability}}

The verification package contains standalone sources, build
instructions, and expected results for the bounded support-eight-to-ten
replay. The analytic multiplicity, rank-phase, and atom-identity
arguments are contained in the manuscript; no separate executable
verification is claimed for them.

\hypertarget{references}{%
\section{References}\label{references}}

\begin{enumerate}
\def\labelenumi{\arabic{enumi}.}
\tightlist
\item
  M. Freeze and W. A. Schmid, ``Remarks on a Generalization of the
  Davenport Constant,'' \emph{Discrete Mathematics} \textbf{310},
  3373-3389 (2010). DOI: 10.1016/j.disc.2010.07.028
\item
  Y. Fan, W. Gao, G. Wang, Q. Zhong, and J. Zhuang, ``On Short Zero-Sum
  Subsequences of Zero-Sum Sequences,'' \emph{Electronic Journal of
  Combinatorics} \textbf{19}(3), P31 (2012). DOI: 10.37236/2602
\item
  K. Zhao, ``On Zero-Sum Subsequences in a Finite Abelian Group of
  Length Not Exceeding a Given Number,'' arXiv:2506.21383v1
  {[}math.CO{]} (2025), Lemma 4.4.
\item
  W. Gao, A. Geroldinger, and W. A. Schmid, ``Inverse Zero-Sum
  Problems,'' \emph{Acta Arithmetica} \textbf{128}, 245-279 (2007). DOI:
  10.4064/aa128-3-5
\item
  Q. Zhong, ``On the Inverse Problem of the \(k\)-th Davenport Constants
  for Groups of Rank 2,'' \emph{Combinatorica} \textbf{45}, article 31
  (2025). DOI: 10.1007/s00493-025-00153-3
\item
  W. D. Gao, Q. H. Hou, W. A. Schmid, and R. Thangadurai, ``On Short
  Zero-Sum Subsequences II,'' \emph{Integers} \textbf{7}, A21 (2007).
\item
  J. E. Olson, ``A Combinatorial Problem on Finite Abelian Groups, I,''
  \emph{Journal of Number Theory} \textbf{1}(1), 8-10 (1969). DOI:
  10.1016/0022-314X(69)90021-3
\end{enumerate}

\end{document}
